\section{Related Work}
\label{sec:related_work}

\subsection{Flow Matching on Euclidean and Riemannian Manifolds}
\label{sec:rel_flow}
Continuous Normalizing Flows (CNFs) parametrise generative transformations between a tractable base distribution and a target data distribution via time-dependent ordinary differential equations (ODEs) \citep{chen2018neural}. Flow Matching (FM) \citep{lipman2023flow, albergo2023building, liu2023flow} broke the historical computational bottleneck of CNFs by directly regressing a neural vector field onto closed-form target paths, enabling simulation-free training. Subsequent advances incorporated minibatch optimal transport couplings \citep{pooladian2023multisample, tong2023improving} to straighten probability trajectories and reduce inference latency.

Recognising that physical and structural data frequently inhabit non-Euclidean spaces, Riemannian Flow Matching (RFM) \citep{chen2024flow} extended vector field regression to general Riemannian manifolds. Recently, the specific polar or cylindrical decoupled manifold construction ($\mathbb{R}^+ \times S^{d-1}$) has been explored in representation learning and abstract feature alignment. For instance, Direct Product Flow Matching (DP-FM) \citep{chen2026dpfm} applies a decoupled metric to Vision-Language Model latent features to resolve angular dynamics distortion during few-shot adaptation, while DisRFM \citep{wang2026disrfm} utilises polar Riemannian flows for structure-preserving graph domain adaptation. However, applying this geometric construction directly to raw, high-dimensional complex wavefields (e.g., $D=4096$) presents unique scaling challenges regarding spatial coherence and optimal transport, which remain largely unexplored.

\subsection{Deep Generative Models for Complex-Valued Signals}
\label{sec:rel_mri}
The processing of complex-valued physical wavefields, such as Magnetic Resonance Imaging (MRI) or Audio STFTs, inherently requires modelling both magnitude and phase \citep{bernstein2004handbook, trabelsi2018deep}. While supervised networks have successfully adopted complex convolutions \citep{cole2021analysis}, modern generative frameworks (such as Score-Based Diffusion Models \citep{song2020score, chung2022score}) almost universally operate strictly on real-valued magnitude images, or naively treat real and imaginary parts as two Cartesian channels in $\mathbb{R}^2$, inheriting the metric of the plane.

Very recently, a few works have begun exploring generative models for complex signals. \citet{schlimbach2026complex} combine a complex-valued variational autoencoder with latent flow matching for brain MRI, and model the phase in a learned Euclidean latent space rather than on its circular domain. Other works, such as PhaseGen \citep{rempe2025phasegen} for MRI or UniverSR \citep{choi2025universr} for audio, similarly default to Euclidean formulations. As we demonstrate in this paper, Cartesian interpolation structurally induces origin singularities and extreme angular velocities. 

\subsection{Optimal Transport Traps in Flow Matching}
\label{sec:rel_ot}
Minibatch optimal transport (OT) is commonly used to pair noise and data in flow matching \citep{pooladian2023multisample, tong2023improving}. Related couplings under a Riemannian cost have been used on other manifolds: FoldFlow-OT \citep{bose2024foldflow} for protein backbones on $\mathrm{SE}(3)^N$, and Metric Flow Matching \citep{kapusniak2024metric} with geodesics of a learned, data-induced metric. Three aspects distinguish our setting from prior work: the modality consists of complex-valued fields; the manifold is $([0, \infty) \times S^1)^{H \times W}$, whose phase factor is a circle at every pixel; the research question: how the coupling behaves as the field dimension grows and why it must not be factorised. The population-level coupling induced by minibatch OT, as well as its dependence on batch size, has recently been characterised by \citet{boite2026minibatch}. Because minibatch OT scales at least quadratically with the batch, cheaper couplings have been proposed: One-Sided Quantile Coupling (QC-FM, \citealp{kim2026onesided}) constructs each source sample from one-dimensional rank statistics of the data along random projections, leaving every data sample intact. A different shortcut, which we show must be avoided, factorises the coupling itself: matching coordinates (amplitude and phase) or spatial patches with independent permutations. This reassembles data samples out of parts of different samples and destroys the joint data distribution (Section~\ref{sec:experiments}).
